% !TEX root = ../../../main.tex
% Block card: catalogue.shape_families -- "Shape families" (Part I, the blocks chapter).
\cardfield{What it does}
Sorts windows into families of similar shape with the Library's own methods, leaves out the
windows that fit no family, and says for each family whether it stands out from noise.

\cardfield{When to use it}
After \blockref{catalogue.event_windows} (or \blockref{preprocessing.trace_shape}), when you want
to know which event shapes recur, by the same computation the Library's grouping editor runs, so
a family here is a family there.

\begin{plainterms}[Families, and what is ``omitted'']
Sorting shapes is like sorting a box of buttons by look. Most methods would force every button
into some pile, even the odd ones out. Here a window that fits no family is \emph{omitted}: set
aside with a reason, rather than forced in. Reasons include ``not in a dense family'' (a
straggler), ``past cut'' (too far from every family's typical member) and ``group too small''
(its pile had fewer members than the minimum, ten by default). Noise also forms piles, so each
family is compared with noise of the recording's own kind, and flagged when it is no better.
\end{plainterms}

\cardfield{What goes in, what comes out}
In go windows. Out comes a family number for each window (1 is the largest family), or $-1$ with
the reason it was left out. The block's page shows each family's typical member, its noise flag
and how alike its members are.

\cardfield{Settings that matter}
\emph{method}: \texttt{hdbscan} (the default) finds crowds of similar shapes and chooses the
number of families itself; tune \emph{hdbscan\_min\_cluster\_size}. \texttt{kshape} makes
\emph{kshape\_k} families, comparing shapes after sliding one along the other; 0 lets it choose.
\texttt{ward} cuts a tree at \emph{ward\_cut}, places every window and tests nothing against
noise. \emph{min\_group} is the smallest family kept, for every method.

\cardfield{Watch out for}
Size does not enter: every window is resampled and z-normalised, so a small wiggle and a large
drop of the same profile land together. This block ignores the train/test roles: every window is
grouped, test windows included. For a training-only clustering of a pool use
\blockref{catalogue.shape_cluster}.
